%=====================================================================
\chapter{Decision Tree Learning}\label{ch:trees}
%=====================================================================
\locator{3}{Part III --- Supervised Learning Algorithms}

\begin{outcomes}
By the end of this chapter you will be able to:
\begin{tight}
  \item \textbf{Classify} an example with a decision tree, and \textbf{read} a
        tree as a set of rules.
  \item \textbf{Compute} entropy and information gain, and \textbf{trace} ID3's
        choice of attribute at each node by hand.
  \item \textbf{Handle} continuous attributes, attributes with many values
        (gain ratio) and the Gini index.
  \item \textbf{Explain} why trees overfit noisy data, and \textbf{apply}
        reduced-error, rule post-pruning and cost-complexity pruning.
  \item \textbf{State} the inductive bias of ID3.
\end{tight}
\end{outcomes}

\begin{prereq}
\chref{math} (entropy), \chref{concept} (hypothesis spaces and inductive bias)
and \chref{generalization} (overfitting, noise, choosing a setting by
cross-validation).
\end{prereq}

\begin{keyterms}
decision tree $\bullet$ internal node, branch, leaf $\bullet$ ID3 $\bullet$
entropy $\bullet$ information gain $\bullet$ split information $\bullet$ gain
ratio $\bullet$ Gini index $\bullet$ CART $\bullet$ threshold split $\bullet$
overfitting $\bullet$ pre-pruning $\bullet$ reduced-error pruning $\bullet$ rule
post-pruning $\bullet$ cost-complexity pruning $\bullet$ Occam's razor
\end{keyterms}

\section{Trees as Hypotheses}

A \term{decision tree} classifies an example by asking a sequence of questions
about its attributes. Each \term{internal node} tests one attribute, each
\term{branch} is one outcome of the test, and each \term{leaf} gives a class.
To classify, start at the root and follow the branches that match the example
until a leaf is reached.

\dsfig{%
\begin{tikzpicture}[font=\scriptsize,
  attr/.style={rounded corners=3pt, draw=primary, fill=primary!10, text=primary!80!black,
               font=\scriptsize\bfseries, minimum width=2.0cm, minimum height=0.65cm},
  yes/.style={rounded corners=8pt, draw=greenhl, fill=greenhl!16, text=greenhl!40!black,
              font=\scriptsize\bfseries, minimum width=1.2cm, minimum height=0.55cm},
  no/.style={rounded corners=8pt, draw=accent, fill=accent!12, text=accent!80!black,
             font=\scriptsize\bfseries, minimum width=1.2cm, minimum height=0.55cm},
  br/.style={font=\tiny, text=gray!55!black, fill=white, inner sep=1pt},
  e/.style={gray!60, line width=0.9pt}]
\node[attr] (o) at (0,0) {Requirements};
\node[attr] (h) at (-4.0,-1.7) {Load};
\node[yes] (ov) at (0,-1.7) {Yes};
\node[attr] (w) at (4.0,-1.7) {Dependencies};
\node[no]  (h1) at (-5.3,-3.3) {No};
\node[yes] (h2) at (-2.7,-3.3) {Yes};
\node[no]  (w1) at (2.7,-3.3) {No};
\node[yes] (w2) at (5.3,-3.3) {Yes};
\draw[e] (o) -- node[br] {Changing} (h);
\draw[e] (o) -- node[br] {Frozen} (ov);
\draw[e] (o) -- node[br] {Vague} (w);
\draw[e] (h) -- node[br] {High} (h1);
\draw[e] (h) -- node[br] {Normal} (h2);
\draw[e] (w) -- node[br] {Many} (w1);
\draw[e] (w) -- node[br] {Few} (w2);
\node[font=\tiny, text=gray!55!black] at (-5.3,-3.85) {S1, S2, S8};
\node[font=\tiny, text=gray!55!black] at (-2.7,-3.85) {S9, S11};
\node[font=\tiny, text=gray!55!black] at (0,-2.25) {S3, S7, S12, S13};
\node[font=\tiny, text=gray!55!black] at (2.7,-3.85) {S6, S14};
\node[font=\tiny, text=gray!55!black] at (5.3,-3.85) {S4, S5, S10};
\end{tikzpicture}}%
{The tree ID3 learns from the fourteen past sprints. Under each leaf are the training
sprints that reach it; every leaf is pure. It is equivalent to the rule: the goal is met
if (requirements Changing and load Normal) or requirements Frozen or (requirements
Vague and Few dependencies).}{sprint-tree}

A tree is a disjunction of conjunctions: each root-to-leaf path is an AND of tests,
and the tree is the OR of the paths that end in ``Yes''. That makes trees far more
expressive than the conjunctions of \chref{concept} --- a large enough tree can
represent \emph{any} discrete function --- and it makes them readable. Trees
suit problems where examples are attribute--value pairs, the target is discrete,
disjunctions may be needed, and the training data may contain errors or missing
values.

\section{Choosing the Question: Entropy and Information Gain}

Fourteen past sprints, each described by four attributes, and whether the sprint
met its goal. \emph{Requirements} were Changing, Frozen or Vague during the
sprint; \emph{Pressure} is the deadline pressure; \emph{Load} is whether the team
was carrying more work than usual; \emph{Dependencies} counts waits on other
teams. (This is Mitchell's \emph{PlayTennis} table with the attributes renamed,
so every number below can be checked against the set text.)

\rowcolors{2}{white}{lightbg}
\begin{longtable}{@{}L{1.3cm}L{2.5cm}L{2.2cm}L{2.0cm}L{2.6cm}L{2.2cm}@{}}
\toprule
\rowcolor{primary!15}
\textbf{Sprint} & \textbf{Requirements} & \textbf{Pressure} & \textbf{Load} & \textbf{Dependencies} & \textbf{Goal met} \\
\midrule
\endhead
S1 & Changing & High & High & Few & No \\
S2 & Changing & High & High & Many & No \\
S3 & Frozen & High & High & Few & Yes \\
S4 & Vague & Medium & High & Few & Yes \\
S5 & Vague & Low & Normal & Few & Yes \\
S6 & Vague & Low & Normal & Many & No \\
S7 & Frozen & Low & Normal & Many & Yes \\
S8 & Changing & Medium & High & Few & No \\
S9 & Changing & Low & Normal & Few & Yes \\
S10 & Vague & Medium & Normal & Few & Yes \\
S11 & Changing & Medium & Normal & Many & Yes \\
S12 & Frozen & Medium & High & Many & Yes \\
S13 & Frozen & High & Normal & Few & Yes \\
S14 & Vague & Medium & High & Many & No \\
\bottomrule
\end{longtable}

ID3 grows the tree top-down and greedily: at each node, it asks which single
attribute best separates the examples that reach it, puts that attribute there,
and recurses on each branch. ``Best separates'' is measured with the entropy of
\chref{math}.

\begin{definitionbox}[Information Gain]
The \term{information gain} of attribute $A$ on example set $S$ is the expected
reduction in entropy from splitting $S$ on $A$:
\[
\text{Gain}(S, A) = \text{Entropy}(S) - \sum_{v \in \text{Values}(A)}
\frac{|S_v|}{|S|}\,\text{Entropy}(S_v)
\]
where $S_v$ is the subset of $S$ with $A = v$. ID3 puts the attribute with the
highest gain at the node.
\end{definitionbox}

\dsfig{%
\begin{tikzpicture}
\begin{axis}[width=8.4cm, height=4.8cm, axis lines=left,
  xlabel={proportion of positive examples $p$}, ylabel={impurity},
  tick label style={font=\tiny}, label style={font=\scriptsize},
  xmin=0, xmax=1, ymin=0, ymax=1.08,
  legend style={font=\tiny, draw=gray!40, at={(0.5,0.03)}, anchor=south}]
\addplot[primary, line width=1.3pt, domain=0.001:0.999, samples=120]
  {-x*ln(x)/ln(2)-(1-x)*ln(1-x)/ln(2)};
\addlegendentry{entropy (bits)}
\addplot[goldhl, line width=1.2pt, domain=0:1, samples=60] {2*x*(1-x)};
\addlegendentry{Gini index}
\addplot[only marks, mark=*, mark size=1.8pt, accent] coordinates {(0.643,0.940)};
\node[font=\tiny, accent, anchor=north] at (axis cs:0.643,0.9) {9 yes, 5 no: 0.940};
\end{axis}
\end{tikzpicture}}%
{Two impurity measures for a two-class node. Both are zero for a pure node and largest
for an even split; they almost always pick the same attribute.}{impurity}

\begin{worked}[ID3's First Choice on the Sprint Data]
\textbf{The root.} 9 Yes, 5 No: $\text{Entropy}(S) = 0.940$ (computed in \chref{math}).

\smallskip
\textbf{Requirements} splits $S$ into Changing (2 Yes, 3 No: entropy 0.971), Frozen
(4, 0: entropy 0) and Vague (3, 2: entropy 0.971):
\[
\text{Gain}(S, \text{Requirements}) = 0.940 - \tfrac{5}{14}(0.971) - \tfrac{4}{14}(0)
- \tfrac{5}{14}(0.971) = \mathbf{0.246}
\]
\textbf{Load}: High (3, 4: 0.985), Normal (6, 1: 0.592):
\[
\text{Gain}(S, \text{Load}) = 0.940 - \tfrac{7}{14}(0.985) - \tfrac{7}{14}(0.592) = \mathbf{0.151}
\]
\textbf{Dependencies}: Few (6, 2: 0.811), Many (3, 3: 1.000):
\[
\text{Gain}(S, \text{Dependencies}) = 0.940 - \tfrac{8}{14}(0.811) - \tfrac{6}{14}(1.000) = \mathbf{0.048}
\]
\textbf{Pressure}: $\text{Gain} = \mathbf{0.029}$ (check it yourself).

\smallskip
\textbf{Requirements wins} and goes at the root. Its Frozen branch is already pure,
so it becomes a Yes leaf.

\smallskip
\textbf{The Changing branch} (S1, S2, S8, S9, S11; entropy 0.971) is split the
same way, using only those five sprints: Gain(Load) $= 0.971$, Gain(Pressure)
$= 0.571$, Gain(Dependencies) $= 0.020$. Load separates them perfectly. On the
Vague branch Dependencies does the same. Every leaf is now pure, and ID3 stops
with the tree of \dsref{sprint-tree}.
\end{worked}

\section{Beyond the Basic Algorithm}

\textbf{Continuous attributes.} A numeric attribute is turned into a yes/no test
$A \le t$. The only thresholds worth trying are midpoints between adjacent sorted
values where the class changes; each is scored by information gain like any
other attribute.

\begin{worked}[Choosing a Threshold]
Six sprints in which $40, 48, 60, 72, 80, 90$ per cent of the stories had been
refined (given acceptance criteria) before the sprint began, and the goal was met
No, No, Yes, Yes, Yes, No. The class changes between 48 and 60 and between 80 and
90, so the candidates are $t = 54$ and $t = 85$. The set has 3 Yes and 3 No:
entropy 1.
\begin{tight}
  \item $t = 54$: $\le 54$ holds \{No, No\} (entropy 0); $> 54$ holds 3 Yes, 1 No
        (entropy 0.811). Gain $= 1 - \tfrac{4}{6}(0.811) = \mathbf{0.459}$.
  \item $t = 85$: $\le 85$ holds 3 Yes, 2 No (0.971); $> 85$ holds \{No\} (0).
        Gain $= 1 - \tfrac{5}{6}(0.971) = \mathbf{0.191}$.
\end{tight}
The test Refined $\le 54$ is chosen.
\end{worked}

\textbf{Attributes with many values.} Information gain favours attributes with
many values. An attribute such as \emph{SprintID} splits the training sprints into
fourteen pure singletons --- maximal gain, zero use. C4.5 corrects this with the
\term{gain ratio}, dividing the gain by the \term{split information}, the entropy
of the split itself:
\[
\text{SplitInfo}(S, A) = -\sum_{v} \frac{|S_v|}{|S|}\log_2 \frac{|S_v|}{|S|},
\qquad
\text{GainRatio}(S, A) = \frac{\text{Gain}(S, A)}{\text{SplitInfo}(S, A)}
\]
For Requirements (5, 4, 5 sprints) SplitInfo $= 1.577$ and the gain ratio is
$0.246/1.577 = 0.156$. For SprintID, SplitInfo $= \log_2 14 = 3.81$, which heavily
discounts its gain.

\textbf{The Gini index.} CART, and scikit-learn by default, uses
$\text{Gini}(S) = 1 - \sum_i p_i^2$ instead of entropy. At the root of the sprint
data it is $1 - (9/14)^2 - (5/14)^2 = 0.459$, falling to 0.343 after splitting on
Requirements. \dsref{impurity} shows why the choice rarely matters: the two
curves have the same shape.

\textbf{Missing values.} An example missing the tested attribute can be sent down
the most common branch, or split fractionally among all branches in proportion to
their training frequencies (C4.5's approach).

\section{Overfitting, Noisy Data and Pruning}

ID3 keeps splitting until every leaf is pure. With noisy data that means growing
branches whose only purpose is to fit the mislabelled examples: add one wrongly
recorded sprint --- (Changing, High, Normal, Many, No) --- and ID3 must grow a new
subtree under Changing/Normal just to isolate it. \dsref{tree-overfit} shows the
effect on a dataset in which 15\% of the labels were flipped.

\dsfig{%
\begin{tikzpicture}
\begin{axis}[width=9.6cm, height=5.0cm, axis lines=left,
  xlabel={size of tree (number of nodes)}, ylabel={accuracy},
  tick label style={font=\tiny}, label style={font=\scriptsize},
  xmin=0, xmax=235, ymin=0.68, ymax=1.02,
  legend style={font=\tiny, draw=gray!40, at={(0.98,0.5)}, anchor=east}]
\addplot[primary, line width=1.2pt, mark=*, mark size=1.2pt] coordinates
  {(3,0.773) (7,0.792) (11,0.809) (15,0.821) (23,0.841) (31,0.851) (47,0.873) (63,0.893) (95,0.924) (127,0.951) (191,0.988) (225,1.000)};
\addlegendentry{on training data}
\addplot[accent, line width=1.2pt, mark=square*, mark size=1.2pt] coordinates
  {(3,0.724) (7,0.749) (11,0.773) (15,0.780) (23,0.792) (31,0.792) (47,0.779) (63,0.767) (95,0.776) (127,0.761) (191,0.747) (225,0.744)};
\addlegendentry{on test data}
\draw[greenhl!70!black, dashed, line width=0.9pt] (axis cs:27,0.68) -- (axis cs:27,0.792);
\node[font=\tiny, greenhl!45!black, anchor=west] at (axis cs:31,0.705)
     {best test accuracy at about 25--30 nodes};
\end{axis}
\end{tikzpicture}}%
{Accuracy against tree size on data with 15\% label noise. Training accuracy climbs to
100\% as the tree memorises the noise; test accuracy peaks early and then declines.}{tree-overfit}

\begin{definitionbox}[Pruning]
\term{Pre-pruning} stops growth early: a maximum depth, a minimum number of
examples per leaf, or a minimum gain. \term{Post-pruning} grows the full tree and
then cuts it back:
\begin{tight}
  \item \term{Reduced-error pruning}: using a separate validation set, repeatedly
        replace the subtree whose removal most improves validation accuracy by a
        leaf labelled with its majority class; stop when every removal hurts.
  \item \term{Rule post-pruning} (C4.5): convert the tree into one rule per
        root-to-leaf path, drop any precondition whose removal does not hurt the
        rule's estimated accuracy, and sort the rules by accuracy.
  \item \term{Cost-complexity pruning} (CART): minimise
        $\text{error} + \alpha \times (\text{number of leaves})$; larger $\alpha$ gives
        smaller trees, and $\alpha$ is chosen by cross-validation.
\end{tight}
\end{definitionbox}

\begin{conceptbox}[Why Rules Are Easier to Prune Than Trees]
Pruning a tree removes a node and everything below it --- a decision about the
root's test affects every path through it. Converting to rules first lets each
path be simplified on its own: the test Requirements = Changing can be dropped
from one rule and kept in another. It also produces a list of plain if--then
statements, which is what a project manager actually wants to read --- and can
turn straight into a risk checklist.
\end{conceptbox}

\section{The Inductive Bias of ID3}

ID3 searches a \emph{complete} hypothesis space --- every discrete function is
some tree --- so, unlike \textsc{Candidate-Elimination}, it does not restrict
what can be learned. Its bias lies in its \emph{search}: it is greedy, never
backtracks, and so prefers \textbf{shorter trees} and trees that put
\textbf{high-gain attributes near the root}. That is a \term{preference bias}
(an ordering over hypotheses), where \textsc{Candidate-Elimination}'s was a
\term{restriction bias} (a limit on which hypotheses exist). Its justification is
a form of \term{Occam's razor}: there are far fewer short trees than long ones,
so a short tree that fits the data is less likely to fit it by coincidence.

\begin{pitfall}
\begin{tight}
  \item \textbf{Letting the tree grow until the leaves are pure} on noisy data. It
        will fit the noise; prune, or cap the depth.
  \item \textbf{Keeping an ID-like attribute.} It has maximal information gain and
        no predictive value. Drop it, or use the gain ratio.
  \item \textbf{Reading the root attribute as ``the most important feature''.}
        Greedy splitting on correlated features picks one almost arbitrarily.
  \item \textbf{Expecting stability.} A small change in the data can change the
        root and hence the whole tree --- the high variance that
        \chref{ensembles} exploits.
  \item \textbf{Scaling features for a tree.} Unnecessary: a threshold test is
        unchanged by any monotone rescaling.
\end{tight}
\end{pitfall}

%---------------------------------------------------------------------
\begin{lab}[ID3 by Hand, Then Pruning by Cross-Validation]
\begin{lstlisting}[language=Python]
import numpy as np
import pandas as pd
from sklearn.datasets import make_classification
from sklearn.model_selection import train_test_split, cross_val_score
from sklearn.tree import DecisionTreeClassifier, export_text

# --- 1. ID3's first choice, computed by hand -----------------------------
rows = ["Changing High High Few No", "Changing High High Many No",
        "Frozen High High Few Yes", "Vague Medium High Few Yes",
        "Vague Low Normal Few Yes", "Vague Low Normal Many No",
        "Frozen Low Normal Many Yes", "Changing Medium High Few No",
        "Changing Low Normal Few Yes", "Vague Medium Normal Few Yes",
        "Changing Medium Normal Many Yes", "Frozen Medium High Many Yes",
        "Frozen High Normal Few Yes", "Vague Medium High Many No"]
df = pd.DataFrame([r.split() for r in rows],
                  columns=["Reqs", "Pressure", "Load", "Deps", "Met"])


def entropy(labels):
    p = labels.value_counts(normalize=True)
    return float(-(p * np.log2(p)).sum())


def gain(data, attr):
    rest = sum(len(g) / len(data) * entropy(g.Met)
               for _, g in data.groupby(attr))
    return entropy(data.Met) - rest


for a in ["Reqs", "Pressure", "Load", "Deps"]:
    print(f"Gain(S, {a:8s}) = {gain(df, a):.3f}")
changing = df[df.Reqs == "Changing"]
print("inside Changing:", {a: round(gain(changing, a), 3)
                           for a in ["Pressure", "Load", "Deps"]})

# --- 2. noisy data: a tree grown without limit, then pruned -------------
X, y = make_classification(n_samples=1500, n_features=10, n_informative=4,
                           flip_y=0.15, random_state=4)  # 15% labels flipped
X_tr, X_te, y_tr, y_te = train_test_split(X, y, test_size=0.5,
                                          random_state=0)
full = DecisionTreeClassifier(criterion="entropy", random_state=0)
full.fit(X_tr, y_tr)
print(f"unpruned: {full.tree_.node_count:3d} nodes, "
      f"train {full.score(X_tr, y_tr):.3f}, "
      f"test {full.score(X_te, y_te):.3f}")

# cost-complexity pruning: choose alpha by 5-fold CV on the training data
alphas = full.cost_complexity_pruning_path(X_tr, y_tr).ccp_alphas[:-1]
cv = [cross_val_score(DecisionTreeClassifier(criterion="entropy",
      ccp_alpha=a, random_state=0), X_tr, y_tr, cv=5).mean() for a in alphas]
best = alphas[int(np.argmax(cv))]
pruned = DecisionTreeClassifier(criterion="entropy", ccp_alpha=best,
                                random_state=0).fit(X_tr, y_tr)
print(f"pruned:   {pruned.tree_.node_count:3d} nodes, "
      f"train {pruned.score(X_tr, y_tr):.3f}, "
      f"test {pruned.score(X_te, y_te):.3f}  (alpha = {best:.4f})")
print(export_text(pruned, max_depth=1))
\end{lstlisting}
Part 1 prints 0.247 for Requirements and 0.152 for Load where the worked example
has 0.246 and 0.151: the code does not round the intermediate entropies. Part 2 stands in for 1{,}500 past
sprints described by ten numeric features, 15\% of them recorded with the wrong
outcome. Pruning cuts the tree from 207 to 63 nodes and \emph{raises} test
accuracy from 0.72 to 0.75 while lowering training accuracy --- the signature of
removing overfitting.
\end{lab}

%---------------------------------------------------------------------
\begin{checkpoint}
\begin{tightnum}
  \item A node holds 8 positive and 8 negative examples. A test splits it into
        (6+, 2$-$) and (2+, 6$-$). Compute the information gain.
  \item Why does information gain favour an attribute such as \emph{TaskID}?
        What fixes it?
  \item What is the difference between a restriction bias and a preference bias?
        Which does ID3 have?
\end{tightnum}
\end{checkpoint}

%---------------------------------------------------------------------
\begin{chaptersummary}
\begin{tight}
  \item A decision tree is a disjunction of conjunctions of attribute tests; it
        can represent any discrete function and reads as a set of rules.
  \item ID3 grows the tree greedily, choosing at each node the attribute with the
        greatest information gain; on the sprint data, Requirements (0.246) wins
        at the root.
  \item Numeric attributes become threshold tests at class boundaries; the gain
        ratio corrects the preference for many-valued attributes; Gini behaves
        like entropy.
  \item Grown to purity on noisy data, trees overfit. Pre-prune with depth or
        leaf-size limits, or post-prune by reduced error, rules or cost
        complexity, choosing the strength by validation.
  \item ID3's inductive bias is a preference for short trees with high-gain
        attributes near the root --- Occam's razor.
\end{tight}
\end{chaptersummary}

\begin{reviewq}
  \item \textbf{(MCQ)} The entropy of a node with 4 positive and 4 negative
        examples is: (a)~0 (b)~0.5 (c)~1 (d)~2.
  \item \textbf{(MCQ)} C4.5's gain ratio divides information gain by: (a)~the
        number of examples (b)~the split information (c)~the Gini index (d)~the
        tree depth.
  \item \textbf{(MCQ)} Cost-complexity pruning with a larger $\alpha$ gives:
        (a)~a larger tree (b)~a smaller tree (c)~the same tree (d)~a deeper tree.
  \item \textbf{(Short)} Explain why a tree that classifies every training example
        correctly may be a poor model.
  \item \textbf{(Short)} Describe reduced-error pruning and state what data it
        requires.
  \item \textbf{(Short)} Why do decision trees not need feature scaling?
  \item \textbf{(Applied)} Compute Gain(S, Pressure) for the full sprint data,
        showing the entropy of each of the three branches.
  \item \textbf{(Applied)} On the Vague branch of the sprint data (S4, S5, S6,
        S10, S14), compute the gain of Load, Pressure and Dependencies, and
        confirm that ID3 chooses Dependencies.
\end{reviewq}
